Au cours de la deuxième années, nous avons implanté le rejeu à l'identique avec des événements de fusion. %
Dans le cas de la fusion, il arrive souvent que l'orbite d'origine est supprimé et il n'est alors pas possible de retrouver les orbites fusionnées. %
Nous complétons donc, grâce à la règle, l'orbite d'origine avec un chemin entre
l'orbite d'origine et les orbites fusionnées par la règle. %

Nous avons également raffiné nos conditions de détection d'événements pour
présenter une méthode statique de détection d'événement. %
Ces travaux ont été présenté aux Journées du Groupe de Travail Modélisation
Géométrique (GTMG) du 15 et 16 mars 2023. %
Un article a été soumis, toujours pour ces travaux, à la conférence Shape
Modeling International 2023. %
Celui-ci ayant été refusé, nous retravaillons ses motivations en vue de le
soumettre conjointement avec un deuxième article sur le rejeu à la conférence
Computer Graphics \& Visual Computing 2023. %


%%% Local Variables:
%%% mode: latex
%%% TeX-master: "../main"
%%% End:
